\chapter{Conclusiones}

El desarrollo de \textbf{iComanda} permitió construir una solución funcional para
modernizar la operación de restaurantes pequeños y medianos. El proyecto logró
centralizar procesos que normalmente se manejan de forma manual, como la toma de
pedidos, el seguimiento de mesas, la comunicación con cocina, la confirmación de
pagos, el registro de gastos, el control básico de inventario y la generación de
reportes administrativos.

La arquitectura basada en \textbf{Supabase como Backend as a Service} resultó adecuada
para el alcance del proyecto, ya que permitió utilizar PostgreSQL, Auth, Storage,
Realtime, políticas RLS y funciones RPC sin desarrollar un servidor propio. Esta
decisión redujo la complejidad de infraestructura y permitió enfocar el esfuerzo en
la experiencia móvil, las reglas de negocio y la integración entre roles.

La aplicación implementa tres roles internos principales: \texttt{admin},
\texttt{mesero} y \texttt{chef}. El administrador gestiona la configuración del
restaurante, usuarios, mesas, menú, pagos, gastos, inventario y reportes. El mesero
opera pedidos y mesas durante el servicio. El chef visualiza pedidos activos, marca
órdenes listas y controla platos agotados. Además, los clientes pueden acceder a la
carta digital y a la sesión de mesa mediante QR, sin instalar una aplicación.

El trabajo con Scrum en \textbf{3 sprints de 2 semanas} ayudó a ordenar el desarrollo
en incrementos: primero la base técnica y autenticación, luego la operación principal
de pedidos y administración, y finalmente los módulos complementarios de tiempo real,
cocina, inventario, carta digital y reportes. Esta división permitió avanzar de forma
iterativa y cerrar una versión funcional para la entrega final.

\section{Alcance Logrado}

\begin{itemize}[leftmargin=1.5em]
  \item Registro y seguimiento de órdenes con estados \texttt{activa}, \texttt{lista},
        \texttt{entregada} y \texttt{cancelada}.
  \item Gestión administrativa de usuarios, restaurante, mesas y menú.
  \item Flujo de cocina con rol \texttt{chef} y marcado de platos agotados.
  \item Carta digital pública y sesión de cliente por mesa mediante QR.
  \item Llamada del cliente al mesero mediante registros en \texttt{table\_calls}.
  \item Confirmación de pagos QR y registro de fecha efectiva de pago.
  \item Registro de gastos, reportes financieros, ranking de platos y ventas por mesero.
  \item Módulo de inventario con materiales, reposiciones, consumos y stock bajo.
\end{itemize}

\section{Limitaciones Identificadas}

\begin{itemize}[leftmargin=1.5em]
  \item La verificación de pagos QR todavía es manual; no existe integración con una
        pasarela bancaria o billetera digital.
  \item El cliente puede consultar menú y llamar al mesero, pero no crea órdenes
        directamente desde la mesa.
  \item El sistema depende de conectividad para operar en tiempo real; no se implementó
        un modo offline completo.
  \item Las pruebas realizadas se plantean principalmente como manuales y exploratorias;
        queda pendiente incorporar pruebas automatizadas.
  \item El uso del plan gratuito de Supabase puede generar pausas por inactividad y
        límites de recursos si se despliega en producción real.
\end{itemize}

\section{Recomendaciones}

\begin{itemize}[leftmargin=1.5em]
  \item Integrar pagos automáticos con una pasarela compatible con el contexto local.
  \item Añadir pruebas automatizadas con Jest y pruebas end-to-end para flujos críticos.
  \item Implementar caché local para permitir operación parcial sin conexión.
  \item Mejorar el módulo de inventario con alertas, historial avanzado y proyección
        de consumo.
  \item Preparar un despliegue productivo con monitoreo de errores, backups y un plan
        de Supabase que evite pausas automáticas.
  \item Mantener actualizada la documentación del esquema de base de datos, políticas
        RLS y decisiones de arquitectura conforme evolucione el sistema.
\end{itemize}
